\section{Introduction}
\label{sec:intro}

Managers cannot attend to everything. A central question in the theory of the firm is which decisions a manager should keep and which she should hand to subordinates \citep{aghion1997formal, dessein2002authority}. This paper asks how the answer depends on the manager's attention capacity.

We build a simple model in which a manager allocates a fixed stock of attention across tasks. Each task can be decided by the manager herself, at an attention cost, or delegated to a subordinate whose preferences are partially misaligned. Delegation saves attention but loses control. Our main result is that the optimal number of delegated tasks is non-monotone in attention: at low capacity the manager delegates almost everything, at intermediate capacity she keeps the most important tasks, and at high capacity she delegates again because the marginal value of her attention on the remaining tasks falls.

The paper contributes to the literature on authority and attention \citep{garicano2000hierarchies}. Section~\ref{sec:model} presents the model. Section~\ref{sec:results} derives the main results. Section~\ref{sec:conclusion} concludes.
